\begin{table}[t]
\centering
\small
\caption{표 7. 평가 환경 비교}
\label{tab:table7_environments}
\begin{tabular}{lrrrr}
\toprule
항목 & 이전 PC (cu129) & 새 머신 cu129 & 새 머신 cu130 & 팀원 (cu130) \\
\midrule
GPU/드라이버 & RTX 4090, 560.94 & RTX 4090, 595.99.02 & RTX 4090, 595.99.02 & RTX 4090, 616.92 \\
OS & Linux-6.6.87.2-microsoft-standard-WSL2-x86\_64 & Linux-7.0.0-42-generic-x86\_64 & Linux-7.0.0-42-generic-x86\_64 & Linux-6.18.33.2-microsoft-standard-WSL2-x86\_64 \\
torch & 2.13.0+cu129 & 2.13.0+cu129 & 2.13.0+cu130 & 2.13.0+cu130 \\
vllm & 0.27.1+cu129 (GitHub 휠) & 0.27.1+cu129 (GitHub 휠) & 0.27.1 & 0.27.1 \\
transformers & 5.19.0 & 5.19.0 & 5.22.2 & 5.22.2 \\
datasets & 5.1.0 & 5.1.0 & 5.2.0 & 5.2.0 \\
trl & 1.16.0 & 1.16.0 & 1.16.0 & 1.16.0 \\
peft & 0.21.2 & 0.21.2 & 0.21.2 & 0.21.2 \\
torchcodec & 0.17.0+cpu & 0.17.0+cpu & 0.17.0 & 0.17.0 \\
\bottomrule
\end{tabular}
\\[2pt]{\footnotesize 새 머신 cu130 열은 dpo\_dv1\_re의 cu130 평가 기록. vllm은 cu129에서 GitHub 릴리스 휠(+cu129), cu130에서 PyPI 기본.}
\end{table}
